%%%%%%%%%%%%%%%%%%%%%%%%%%%%%%%%%%%%%%%%%%%%%%%%%
% DOCUMENT SETUP: Do not make any changes to this section.
%%%%%%%%%%%%%%%%%%%%%%%%%%%%%%%%%%%%%%%%%%%%%%%%%

\documentclass[a4paper,10.5pt,twoside]{article}
\hyphenpenalty=8000
\textwidth=125mm
\textheight=200mm

\usepackage{fancyhdr}
\pagestyle{fancy}
\fancyhead{}
\fancyfoot{}
\raggedbottom

\usepackage{xurl}
\usepackage{graphicx}
\usepackage{alltt}
\usepackage{subcaption}
\usepackage{amsmath}
\usepackage[hidelinks]{hyperref}
\urlstyle{same}
\usepackage[T1]{fontenc}
\usepackage[utf8]{inputenc}
\usepackage{lmodern}
\usepackage{csquotes}
\usepackage{comment}
\usepackage{float}
\usepackage{xcolor}
\usepackage[a4paper, top=3cm, bottom=2cm, left=3cm, right=3cm]{geometry}
\usepackage[english]{babel}

\pagenumbering{arabic}
\setcounter{page}{1}

\hypersetup{
    colorlinks=true,
    linkcolor=blue,
    citecolor=blue,
    urlcolor=blue
}

%%%%%%%%%%%%%%%%%%%%%%%%%%%%%%%%%%%%%%%%%%%%%%%%%
\begin{document}

\fancyhead[LE]{\thepage\ \ \ MBD Research Proposal}
\fancyhead[RO]{Multimodal Misinformation Detection\ \ \ \thepage}

\begin{center}
{\Large \textbf{Multimodal Image--Text Misinformation Detection Using Pretrained Contrastive Models}}\\[12pt]
{\normalsize
Ishimwe Karekezi Guy Gael\\
Lynne Chepkwony\\
Emile Lucky Muhigira
}
\end{center}

\section{Project Title and Context}

The spread of fake news on social media increasingly relies on the combination of images with contradictory or irrelevant text in order to mislead audiences. In many cases, the image itself is authentic, while the accompanying text distorts its meaning, a challenge that is difficult to address using systems that analyze text or images in isolation. This problem is particularly pronounced in African media ecosystems, where content verification capacity remains limited and the wide diversity of languages complicates the deployment of traditional automated solutions. This project focuses on detecting fake news by measuring the semantic consistency between image content and textual claims, using pretrained multimodal models that are computationally efficient and more scalable than traditional fact-verification approaches.

\section{Research Problem \& Question(s)}

Misinformation frequently takes advantage of semantic inconsistencies between images and accompanying text to mislead audiences without explicitly fabricating content. Many existing multimodal misinformation detection systems depend on large-scale supervised learning or external knowledge bases, approaches that are often expensive and difficult to deploy in low-resource settings. Recent progress in pretrained multimodal representation learning indicates that the degree of semantic alignment between images and text can be an effective signal for identifying misleading content. However, the performance and reliability of such lightweight approaches in non-Western settings, particularly within African contexts, have not yet been sufficiently explored.

The research questions guiding this study are as follows:
\begin{itemize}
    \item Can semantic inconsistency between images and their accompanying text, quantified using pretrained multimodal embeddings, reliably indicate the risk of misinformation?
    \item How robust is this approach when evaluated on African-context examples compared to standard benchmark datasets?
\end{itemize}


\section{Data Source \& Collection Plan}

This project will rely exclusively on existing publicly available multimodal datasets containing image--text pairs and misinformation labels, including Fakeddit, MM-COVID, and VisualNews. These datasets consist primarily of social media posts and news articles with associated images and captions, predominantly in English. A manageable subset of approximately 1,000--5,000 samples will be selected for experimentation and evaluation. No web scraping or large-scale annotation will be performed at this stage. To assess contextual robustness, a small manually curated test set of 10--30 African-context image--text examples will be created solely for evaluation purposes.

\section{Ethics, Risks, and Limitations}

This work extends prior research on multimodal misinformation detection by focusing on a lightweight, pretrained representation-based approach rather than large-scale supervised or knowledge-intensive systems. Ethical risks include exposure to misleading or potentially harmful content during analysis; however, all data used will originate from publicly available datasets, and no personally identifiable information will be collected or redistributed. The system is designed to flag misinformation risk rather than determine factual truth, reducing the likelihood of over-censorship or misuse. Potential biases may arise from dataset composition, particularly the overrepresentation of Western-centric news and social media content. A key limitation of the approach is that semantic inconsistency alone cannot capture all forms of misinformation.

\section{Expected Outcome}

The expected outcome of this project is a functional prototype that takes an image and its accompanying text as input and produces a similarity score together with a misinformation risk assessment. The project will also deliver an analysis report that examines the effectiveness of image--text semantic inconsistency as a signal for detecting misinformation across multiple benchmark datasets. Qualitative evaluation will be conducted to identify failure cases and contextual limitations, with particular attention to content related to African contexts. The system will be deployed as a lightweight web-based demonstration using free hosting platforms such as Hugging Face Spaces or Streamlit Community Cloud. This work aims to provide a practical and reusable baseline for researchers and practitioners developing scalable multimodal content moderation systems.


\section{References}

\begin{thebibliography}{9}

\bibitem{radford2021clip}
Radford, A., Kim, J. W., Hallacy, C., Ramesh, A., Goh, G., Agarwal, S., Sastry, G., Askell, A., Mishkin, P., Clark, J., et al. (2021).
Learning transferable visual models from natural language supervision.
\textit{Proceedings of the 38th International Conference on Machine Learning}.

\bibitem{khattar2020spotfake}
Khattar, D., Goud, J. S., Gupta, S., and Varma, V. (2020).
SpotFake: A multimodal framework for fake news detection.
\textit{Proceedings of the IEEE International Conference on Big Data}.

\bibitem{liu2021visualnews}
Liu, F., Shi, J., and Liu, Y. (2021).
VisualNews: A large-scale dataset for image--caption analysis in news.
\textit{arXiv preprint arXiv:2101.01790}.

\bibitem{jin2021mmcovid}
Jin, Z., Cao, J., Guo, H., Zhang, Y., and Luo, J. (2021).
Multimodal fusion with recurrent neural networks for rumor detection on microblogs.
\textit{ACM Multimedia}.

\bibitem{li2024kamp}
Li, S., Yuan, J., Chen, Z., Zhao, S., Wang, P., and Huang, X. (2024).
Knowledge-aware multimodal pretraining for fake news detection.
\textit{arXiv preprint arXiv:2401.01234}.


\end{thebibliography}

\end{document}
